\section{模型评价与推广}

\subsection{模型优势}

\begin{enumerate}
  \item 覆盖宽度与重叠率全程使用水平投影口径，可与水平测线间距直接比较，避免了斜坡情形下坡面长度与水平距离混用的问题；
  \item 问题一、二为确定性解析几何模型，参数可解释、结果可复算，并在坡度为零、$\beta=0^\circ$ 或 $180^\circ$ 等特殊情形下退化为平底公式，有明确的口径检验；
  \item 问题三的贪心递推利用等深线平行特征把二维布线降为一维递推，33 条测线的缺口复核表明当前方向下测线数不可再减少，结果可解释；
  \item 问题四采用分带策略并保留网格覆盖仿真，直接对应题目要求的总长度、漏测率与超重叠长度三项指标，并明确区分“指标最优”与“工程折中”两种推荐口径。
\end{enumerate}

\subsection{模型局限与推广}

\begin{enumerate}
  \item 问题一至问题三假设海底为理想平面斜坡、声波沿直线传播，未考虑真实声速剖面、折射和复杂地形；
  \item 问题四仅统计测量线段长度，未计入线段之间的转场、调头和航行连接成本，可能低估实际作业代价；
  \item 分带边界处各带独立布线的覆盖近似会引入少量重叠误差，局部最大相邻重叠率仍达 $49.01\%$；
  \item 问题四未进行随机扰动或全局优化，所得为当前分带与递推规则下的较优方案，而非严格全局最优。
\end{enumerate}

将平面坡度替换为局部梯度估计后，“水平投影覆盖宽度 + 相邻重叠率”的评价框架可迁移到一般起伏海底；若引入声速分层，可在射线追踪中修正声线曲率\cite{ref_book}；若考虑船舶转弯与多船协同，可把分带递推扩展为带转场代价的路径优化问题。
